% =====================================================
% 题型：解三角形辅助角换元选择题 (q_triangle_right_auxiliary.tex)
% =====================================================
% 难度：★★ (中档)
% =====================================================
\newcommand{\qTriangleRightAuxStem}{在 \(\triangle ABC\) 中，角 \(A,B,C\) 的对边分别为 \(a,b,c\)，满足 \(b\cos A + a\sin B = c\)。则角 \(B\) 的值为 \hfill （\quad）}

\newcommand{\qTriangleRightAuxOptions}{%
  \mcoptions[4]{\(\dfrac{\pi}{6}\)}{\(\dfrac{\pi}{4}\)}{\(\dfrac{\pi}{3}\)}{\(\dfrac{\pi}{2}\)}%
}

\newcommand{\qTriangleRightAuxAnswer}{D}

\newcommand{\qTriangleRightAuxSolution}{%
  由正弦定理，可设
  \[
    a=2R\sin A,
    \qquad
    b=2R\sin B,
    \qquad
    c=2R\sin C.
  \]
  代入条件：
  \[
    2R\sin B\cos A
    +2R\sin A\sin B
    =2R\sin C.
  \]
  两边同除以 \(2R\)，得
  \[
    \sin B(\cos A+\sin A)=\sin C.
  \]
  因为 \(C=\pi-A-B\)，所以
  \[
    \sin C=\sin(A+B).
  \]
  因而
  \[
    \sin B(\cos A+\sin A)=\sin(A+B).
  \]
  展开右边：
  \[
    \sin B\cos A+\sin A\sin B
    =\sin A\cos B+\cos A\sin B.
  \]
  两边消去 \(\sin B\cos A\)，得
  \[
    \sin A\sin B=\sin A\cos B.
  \]
  又 \(A\in(0,\pi)\)，故 \(\sin A>0\)，于是
  \[
    \sin B=\cos B.
  \]
  因为 \(B\in(0,\pi)\)，解得
  \[
    B=\frac\pi4.
  \]
  故选 B。%
}

\newcommand{\qTriangleRightAuxDiagnosis}{%
  若能完成边化角但无法继续，多半是没有想到把 \(\sin C\) 写成 \(\sin(A+B)\)。这一步是本题的结构入口。%
}

\newcommand{\qTriangleRightAuxTeachingNotes}{%
  训练“正弦定理边化角 → 和角公式 → 约去必正项”的短链。应提醒学生最后必须结合 \(B\in(0,\pi)\) 确定唯一角。%
}
